\section{Direct stopping for the original Bellman decision}
\label{sec:decisionstopping126}

A continuation critic is useful only if its proposal can be turned into
a qualified economic decision at a competitive complete cost.  The
inherited warm service begins with an already completed conventional
service and then purchases more original information.  The direct
service instead pays for its own actor proposal and its own original
continuous-risk supports, and asks after each support whether the
original decision is already resolved.  It uses the same original
model, continuation policy, candidate feasibility and qualification
tolerances.  Every representation receives this service, including the
ordinary hinge with exactly the neural features and coefficients.
The comparison therefore tests a decision procedure available to all
representations; it does not assign its gains to neural architecture.

The stopping rule separates the action question from the absolute-value
question.  Let $q$ be the original optimal-future action value,
$q^*=\min_{[0,A]}q$, and $J=q+\eta$ the unchanged implemented policy,
where $0\le\eta\le\delta$.  Original derivative intervals and the
proved curvature bounds give a certified optimizer bracket $[\ell,r]$.
At a realized feasible action $a$, write
$\rho(a)=\max\{|a-\ell|,|a-r|\}$ and define the constrained projected
gap
\[
 D_\mu(a;g_-,g_+)=
 \max_{z\in\{g_-,g_+\}}
 \left\{-z d_z-\frac\mu2d_z^2\right\},\qquad
 d_z=\operatorname{proj}_{[-a,A-a]}(-z/\mu),
\]
where $[g_-,g_+]$ encloses the original $q'(a)$.  Retained supports
also give $L_q(a)\le q(a)\le U_q(a)$ and $B\le q^*$.
Proposition~\ref{prop:directcertificate126} proves that
\[
 J(a)-q^*\le G(a):=
 \min\{U_q(a)+\delta-B,\ \delta+D_\mu(a;g_-,g_+),
                              \delta+M\rho(a)\}.
\]
When original endpoint derivative signs certify a stationary optimizer,
$\delta+L\rho(a)^2/2$ is an additional route in this minimum.  The
projected bound also handles capacity optima with nonzero gradients.
No derivative of the arbitrary future-controller error $\eta$ is used.

Qualification requires both $G(a)\le\tau$ and the separate absolute
interval condition $U_q(a)+\delta-L_q(a)\le\omega$, after outward
rounding.  The interval for the original optimum remains a distinct
object: $J^+(a)-G(a)$ is not generally an optimum lower bound.  A
known optimal action can still have an unresolved additive value level,
as Proposition~\ref{prop:separategates126} shows.  This separation identifies action sufficiency independently of the
value refinement still required by the stated output; it does not assert
that the finite implemented acquisition schedule minimizes either cost.  It does not turn optimizer localization into free value
information.

The first completed qualified action needs no comparison to a completed
incumbent.  Every later distinct-action replacement must itself meet
both gates and have a strictly negative upper bound on the original
complete-policy contrast.  In particular,
$\overline\Delta_q(a,b)+\delta<0$ certifies a gain of at least
$-\overline\Delta_q(a,b)-\delta$.  A sharper certificate at the same
action is recorded as requalification, not a policy improvement.

Further action centers use the actually rounded midpoint of the
certified bracket.  After the first actor center, moving to another
center requires the observed contraction
$w_{\rm new}\le w_{\rm old}/2+\epsilon$, where $\epsilon$ is the
actual rounding error relative to the old midpoint.  The finite
precision ladder and resource cap may prevent that condition or the
stronger regret target from being reached.  The service then retains
its ordinarily qualified policy, or pays the complete conventional
fallback if none qualifies.  Lemma~\ref{lem:roundstop126} and
Proposition~\ref{prop:directservice126} prove these claims with the
rounding and payment rules made explicit.  Comparative cost or economic
superiority requires the separately executed matched experiment.
